\documentclass[11pt,reqno]{amsart}

\usepackage[margin=1in]{geometry}
\usepackage{amsmath,amssymb,amsthm,mathtools}
\usepackage{hyperref}
\usepackage{xcolor}
\hypersetup{
  colorlinks=true,
  linkcolor=blue!50!black,
  citecolor=blue!50!black,
  urlcolor=blue!50!black,
  pdfauthor={Vit Planocka},
  pdftitle={The real-rootedness probability of a random cubic polynomial}
}

\theoremstyle{plain}
\newtheorem{theorem}{Theorem}
\newtheorem{lemma}[theorem]{Lemma}
\newtheorem{proposition}[theorem]{Proposition}
\newtheorem{corollary}[theorem]{Corollary}
\theoremstyle{remark}
\newtheorem{remark}[theorem]{Remark}
\newtheorem{example}[theorem]{Example}

\DeclareMathOperator{\vol}{vol}
\newcommand{\RR}{\mathbb{R}}
\newcommand{\Prb}{\mathbb{P}}
\newcommand{\disc}{\Delta}

\title[Real-rootedness of a random cubic]{The real-rootedness probability of a random cubic polynomial: three closed forms}

\author{Vít Plánočka}
\thanks{With assistance from Claude (Anthropic). Machinery, proofs, and their
verification (symbolic, numerical, and partially machine-checked in Lean~4 /
Mathlib) are recorded in the accompanying repository.}

\date{August 2026}

\begin{document}

\begin{abstract}
We compute, in closed form, the probability that a random cubic polynomial has
three real roots, in three related models. For a \emph{monic} cubic
$x^3+ax^2+bx+c$ with $(a,b,c)$ independent and uniform on $[-1,1]$, this
probability is $\tfrac{383}{4860}+\tfrac{\ln 3}{48}$; on $[0,1]$ it is exactly
$\tfrac{1}{2880}$. For the \emph{non-monic} cubic $ax^3+bx^2+cx+d$ with all four
coefficients independent and uniform on $[-1,1]$, it is
$\tfrac{641}{2430}-\tfrac{\ln 3}{24}$. The first two values appear to be new.
The third confirms, with a first proof, a closed form that had circulated
unproved in an informal online forum thread since June~2026. All three
results are verified by independent symbolic and numerical computation, and
by a complete, zero-\texttt{sorry} machine-checked proof of all three
theorems in Lean~4 / Mathlib -- which along the way proves, from scratch, the
classical fact connecting a real cubic's discriminant sign to its real-root
count, seemingly absent from Mathlib until now. We record the reduction
technique for each case -- an elementary band/discriminant-sign argument for
the monic case, and a divergence-theorem argument exploiting the homogeneity
of the discriminant for the non-monic case, which formalization showed can be
replaced by a strictly more elementary slicing argument -- since both appear
to generalize to related open problems, which we list.
\end{abstract}

\maketitle

\tableofcontents

\section{Introduction}

How likely is a ``random'' polynomial to have all of its roots real? For
degree~$2$ this is classical. For a monic quadratic $x^2+bx+c$ with $(b,c)$
uniform on $[-1,1]^2$, real roots occur exactly when $c\le b^2/4$, giving
probability $13/24$; for the full quadratic $ax^2+bx+c$ the answer involves a
logarithm, $\tfrac{41}{72}+\tfrac{\ln 2}{12}$ on $[-1,1]^3$ and
$\tfrac{5}{36}+\tfrac{\ln 2}{6}$ on $[0,1]^3$, the latter apparently first
recorded (without derivation) by J.~D'Aurizio and rederived in
\cite{haldar-chakraborty}. For the analogous \emph{cubic} question the
literature is far thinner. The only paper we could locate that claims exact
results for random cubics with general coefficient distributions is
Li~\cite{li1988}, which is effectively unobtainable (two citations in
thirty-eight years, neither reproducing a value) and, per the recovered
\emph{Mathematical Reviews} notice \cite{li1988-mr}, treats a family of
symmetric-interval distributions that does not include the one-sided case of
our Theorem~\ref{thm:unit-cube} at all, and reaches our Theorem~\ref{thm:sym-cube}
only as an unverifiable special case.

This note settles three instances of the cubic question completely.

\begin{theorem}[Monic cubic, symmetric cube] \label{thm:sym-cube}
Let $(a,b,c)$ be independent and uniform on $[-1,1]$. Then
\[
  \Prb\bigl(x^3+ax^2+bx+c \text{ has three real roots}\bigr)
  = \frac{383}{4860}+\frac{\ln 3}{48}
  = 0.10169434037605886959954086\ldots
\]
\end{theorem}

\begin{theorem}[Monic cubic, unit cube] \label{thm:unit-cube}
Let $(a,b,c)$ be independent and uniform on $[0,1]$. Then
\[
  \Prb\bigl(x^3+ax^2+bx+c \text{ has three real roots}\bigr) = \frac{1}{2880}.
\]
\end{theorem}

\begin{theorem}[Non-monic cubic] \label{thm:nonmonic}
Let $(a,b,c,d)$ be independent and uniform on $[-1,1]$. Then
\[
  \Prb\bigl(ax^3+bx^2+cx+d \text{ has three real roots}\bigr)
  = \frac{641}{2430}-\frac{\ln 3}{24}
  = 0.21801049620261477108898412335868\ldots
\]
\end{theorem}

Theorems~\ref{thm:sym-cube} and \ref{thm:unit-cube} are, to our knowledge, new:
an eight-agent literature sweep together with a targeted search for the exact
constants (Section~\ref{sec:novelty}) found no prior publication of either
value, digit string, or exact expression. Theorem~\ref{thm:nonmonic} is more
subtle. The exact value $\tfrac{641}{2430}-\tfrac{\ln3}{24}$ already existed,
unproved, as the terminus of a long, informally written, partly AI-assisted
derivation posted to a Russian mathematics forum (\texttt{dxdy.ru}) in June
2026. Section~\ref{sec:nonmonic} gives the first proof of this value, together
with five independent lines of numerical and symbolic corroboration.

\subsection*{Motivation and context}
The cubic question is the tractable sibling of a harder open problem: for an
$n\times n$ matrix with i.i.d.\ uniform entries, what is the probability that
all $n$ eigenvalues are real? For $n=2$ this is classical
($49/72$ for entries uniform on $[-1,1]$, \cite{hetzel-liew-morrison}), but
already for $n=3$ no closed form or rigorous nontrivial bound is known; the
canonical statement of the question, posted to Mathematics Stack Exchange in
2020, has $153$ upvotes and, as of this writing, zero answers
\cite{mse-eigenvalue-question}. The obstruction is structural: for Gaussian
entries, orthogonal invariance gives Edelman's exact
$2^{-n(n-1)/4}$ \cite{edelman1997}, but the uniform cube has no such symmetry
to exploit. Moving from matrix-entry space to \emph{coefficient} space --
asking instead for the probability that a random \emph{characteristic
polynomial} is real-rooted -- keeps the same underlying object (the volume of
a semialgebraic region cut out by a discriminant) while dropping the ambient
dimension from $9$ (a $3\times3$ matrix) to $3$--$4$ (a cubic's coefficients),
where it becomes tractable by elementary means. This note is the result of
attacking that lower-dimensional sibling.

\subsection*{Machine verification}
Every calculus step in the proofs of Theorems~\ref{thm:sym-cube}--\ref{thm:nonmonic}
is checked twice over in exact arithmetic (once by hand below, once
independently by computer algebra), confirmed by high-precision quadrature
and by direct Monte Carlo sampling, and -- new for a paper at this level of
the literature -- verified by a \emph{complete, zero-\texttt{sorry}
machine-checked proof of all three theorems} in Lean~4 against Mathlib,
described in Section~\ref{sec:lean}. That formalization also produced a
result of independent interest: a from-scratch, calculus-free proof that a
real cubic's discriminant sign determines its real-root count, a classical
fact we could not find already in Mathlib.

\section{Preliminaries} \label{sec:prelim}

Throughout, $\disc$ denotes the discriminant of a cubic. For the general cubic
$ax^3+bx^2+cx+d$,
\begin{equation} \label{eq:disc4}
  \disc(a,b,c,d) = 18abcd - 4b^3d + b^2c^2 - 4ac^3 - 27a^2d^2,
\end{equation}
and classically, for $a\ne0$: $\disc>0$ iff the cubic has three distinct real
roots; $\disc=0$ iff it has a repeated root (all real); $\disc<0$ iff it has
one real root and a complex-conjugate pair. For the monic case $a=1$ we write
$\disc_3(b,c,d)$ for $\disc(1,b,c,d)$; abusing notation slightly to match the
variable names used in Theorems~\ref{thm:sym-cube}--\ref{thm:unit-cube} (whose
monic cubic is $x^3+ax^2+bx+c$, i.e.\ our $(b,c,d)\mapsto(a,b,c)$),
\begin{equation} \label{eq:disc3}
  \disc_3(a,b,c) = 18abc - 4a^3c + a^2b^2 - 4b^3 - 27c^2.
\end{equation}

\section{The monic cubic}

\subsection{The band reduction}

Write $f(x)=x^3+ax^2+bx+c$ and $g(x)=x^3+ax^2+bx=f(x)-c$. For
$b<a^2/3$, put $s=\sqrt{a^2-3b}$; the critical points of $f$ (equivalently of
$g$) are
\[
  x_-=\frac{-a-s}{3}\ (\text{local max}), \qquad x_+=\frac{-a+s}{3}\ (\text{local min}).
\]
$f$ has three real roots iff $f(x_-)\ge 0\ge f(x_+)$, i.e.\ iff $c$ lies in the
band
\[
  c\in[c_{\mathrm{lo}},c_{\mathrm{hi}}], \qquad
  c_{\mathrm{lo}}=-g(x_-),\quad c_{\mathrm{hi}}=-g(x_+),\qquad
  c_{\mathrm{hi}}-c_{\mathrm{lo}}=\frac{4}{27}s^3.
\]
This is not merely an implication but captures the discriminant exactly: as a
polynomial in $c$,
\begin{equation} \label{eq:disc-factor}
  \disc_3(a,b,c) = -27\,(c-c_{\mathrm{lo}})(c-c_{\mathrm{hi}}).
\end{equation}
\eqref{eq:disc-factor} is verified by direct expansion (and, independently, by
\texttt{sympy}). For $b>a^2/3$ the right-hand quadratic in $c$ has no real
roots and $\disc_3<0$ identically (one real root). Hence for any $c$-window
$[w_{\mathrm{lo}},w_{\mathrm{hi}}]$,
\begin{equation} \label{eq:vol-formula}
  \vol = \iint_{b<a^2/3} \operatorname{len}\bigl([c_{\mathrm{lo}},c_{\mathrm{hi}}]\cap[w_{\mathrm{lo}},w_{\mathrm{hi}}]\bigr)\,db\,da.
\end{equation}

\subsection{Proof of Theorem~\ref{thm:sym-cube}}

\begin{lemma}[Never-clipped] \label{lem:never-clipped}
On $D=\{(a,b): a\in[-1,1],\ b\in[-1,a^2/3]\}$, $c_{\mathrm{hi}}\le 1$, with
equality only at $(a,b)=(-1,-1)$.
\end{lemma}

\begin{proof}
Since $\partial g/\partial a$ at fixed $x$ equals $x^2$ and $g'(x_+)=0$ kills
the chain-rule term at the critical point, $\partial c_{\mathrm{hi}}/\partial a
= -x_+^2\le0$; so $c_{\mathrm{hi}}$ is non-increasing in $a$. As
$b\le a^2/3\le 1/3$ for $a\in[-1,1]$, the point $(-1,b)$ is also in $D$ for
every $(a,b)\in D$, so $c_{\mathrm{hi}}(a,b)\le c_{\mathrm{hi}}(-1,b)$. On the
edge $a=-1$: $\partial c_{\mathrm{hi}}/\partial b=-x_+$ with
$x_+=(1+\sqrt{1-3b})/3>0$, so $c_{\mathrm{hi}}$ is strictly decreasing in $b$
there, with maximum at $b=-1$: $c_{\mathrm{hi}}(-1,-1)=1$ exactly.
\end{proof}

The map $(a,b,c)\mapsto(-a,b,-c)$ preserves $\disc_3$ (direct expansion) and
the cube, and sends $c_{\mathrm{hi}}$ to $-c_{\mathrm{lo}}$; hence
$c_{\mathrm{lo}}\ge-1$ on $D$ as well. \textbf{The band never exits the window
$[-1,1]$}: it touches the boundary only at the two corner points
$(\mp1,-1,\pm1)$, of measure zero. Therefore, using
\eqref{eq:vol-formula} with no clipping correction,
\[
  \vol = \int_{-1}^{1}\int_{-1}^{a^2/3}\frac{4}{27}(a^2-3b)^{3/2}\,db\,da
       = \frac{8}{405}\int_{-1}^{1}(a^2+3)^{5/2}\,da
       = \frac{3064}{1215}+\frac{8}{3}\operatorname{arsinh}\!\frac{1}{\sqrt3}.
\]
Since $\operatorname{arsinh}(1/\sqrt3)=\ln(1/\sqrt3+2/\sqrt3)=\tfrac12\ln 3$,
\[
  \Prb = \vol/8 = \frac{383}{4860}+\frac{\ln 3}{48}. \qedhere
\]

\begin{remark}
The corner-touching geometry is exactly why this is solvable in closed form:
the naive ``main term plus clipping corrections'' program has \emph{all}
corrections identically zero, so the trivial upper bound already equals the
answer. High-precision quadrature landing on that upper bound to sixteen
digits was, in practice, the signal that this was happening.
\end{remark}

\subsection{Proof of Theorem~\ref{thm:unit-cube}}

Now the domain is $a\in[0,1]$, $b\in[0,a^2/3]$, window $[0,1]$.
Since $b\ge0$ forces $s\le a$, both critical points satisfy $x_-\le0\le$
(the argument of the local max) and $x_+\le0$; and $c_{\mathrm{hi}}\ge0$
because $g(0)=0$, $g'(0)=b\ge0$, and $x_+\le0$ is the local minimum, so
$g(x_+)\le0$. A short computation shows $c_{\mathrm{hi}}\le1/27$ throughout the
domain, so the top of the window $[0,1]$ never binds. The bottom does:
$\partial c_{\mathrm{lo}}/\partial b=-x_-\ge0$, and $c_{\mathrm{lo}}=0$ exactly
on the curve $b=a^2/4$ (where $x_-=-a/2$ and $g(-a/2)=0$); so
$c_{\mathrm{lo}}\le0$ for $b\le a^2/4$ and $c_{\mathrm{lo}}\ge0$ for
$a^2/4\le b\le a^2/3$ -- a single sign change. Hence the overlap length is
$c_{\mathrm{hi}}$ on $[0,a^2/4]$ and the full band width on
$[a^2/4,a^2/3]$, and, remarkably, all fractional powers cancel:
\[
  \int_0^{a^2/4}c_{\mathrm{hi}}\,db+\int_{a^2/4}^{a^2/3}\frac{4}{27}(a^2-3b)^{3/2}\,db = \frac{a^5}{480}.
\]
Integrating over $a\in[0,1]$,
\[
  \Prb = \int_0^1\frac{a^5}{480}\,da = \frac{1}{2880}. \qedhere
\]

\section{The non-monic cubic} \label{sec:nonmonic}

Dropping the monic restriction -- letting the leading coefficient $a$ range
uniformly over $[-1,1]$, including near $0$ -- makes the problem genuinely
four-dimensional, and Lemma~\ref{lem:never-clipped}'s trick does not
generalize directly: it is special to a symmetric window of half-width
\emph{exactly} $1$ around a \emph{fixed} leading coefficient.

Writing $V(t)$ for the volume of the real-rooted region inside the cube
$[-t,t]^3$ in the monic-cubic coefficient space (so $V(1)=8\cdot\Prb$ of
Theorem~\ref{thm:sym-cube}), the natural attack conditions on the leading
coefficient and rescales:
\begin{equation} \label{eq:naive-integral}
  \Prb(\text{non-monic}) = \int_0^1 \frac{a^3}{8}\,V(1/a)\,da.
\end{equation}
This requires $V(t)$ for every $t\ge1$, and \eqref{eq:naive-integral} turns
out to be a dead end: high-precision numerics for $p(a):=a^3V(1/a)/8$
(computed to $40$ digits; e.g.\ $p(1/2)=0.1732556347784\ldots$ and
$p(1/3)=0.2277312312100\ldots$) show no recognizable closed form under a
blind PSLQ search against $\{1,\ln2,\ln3,\ln5,\sqrt2,\sqrt3,\pi\}$. $V(1)$ is
elementary only because $t=1$ is exactly where the corner-touching geometry of
Lemma~\ref{lem:never-clipped} applies -- not because $V(t)$ is nice in
general.

\subsection{The discriminant region is a cone}

The route that works is structural rather than a direct generalization of the
band calculus. The discriminant \eqref{eq:disc4} is homogeneous of degree $4$:
$\disc(\lambda x)=\lambda^4\disc(x)$ for $x=(a,b,c,d)$, so the real-rooted set
$R=\{\disc>0\}$ is invariant under $x\mapsto\lambda x$ ($\lambda>0$) --
a cone. Apply the divergence theorem to $F=x/4$ (so $\operatorname{div}F=1$)
on $R\cap[-1,1]^4$. Euler's identity gives $x\cdot\nabla\disc=4\disc$, which
vanishes identically on the lateral boundary $\{\disc=0\}$; so that part of
the surface integral is zero, and only the eight faces of the cube
contribute:
\[
  \vol_4(R\cap[-1,1]^4) = \tfrac14\oint_{\partial(R\cap[-1,1]^4)} F\cdot n\,dS
  = \tfrac12(S_a+S_b+S_c+S_d),
\]
where $S_a,S_b,S_c,S_d$ are the (signed, outward) contributions of the four
faces $a=\pm1,\ldots$ paired by central symmetry
$x\mapsto-x$ (which preserves $\disc$ and the cube). The coefficient reversal
$(a,b,c,d)\mapsto(d,c,b,a)$ -- sending each root $x$ to its reciprocal $1/x$ --
also preserves both $R$ and the cube, and identifies $S_a=S_d$, $S_b=S_c$.
Hence
\begin{equation} \label{eq:cone-reduction}
  \vol_4(R\cap[-1,1]^4)=V(1)+S_b, \qquad \Prb(\text{non-monic})=\frac{V(1)+S_b}{16},
\end{equation}
where $V(1)=S_a=\tfrac{766}{1215}+\tfrac{\ln3}{6}$ (eight times
Theorem~\ref{thm:sym-cube}) and
\[
  S_b = \vol_3\bigl\{(a,c,d)\in[-1,1]^3 : ax^3+x^2+cx+d \text{ has three real roots}\bigr\}
\]
is the volume on the face where the \emph{second} coefficient is pinned to
$1$ -- a genuinely new, but purely three-dimensional, problem, with neither a
$t\to\infty$ tail nor an $a\to0$ endpoint singularity to contend with.

\begin{remark}
We verified the reduction step \eqref{eq:cone-reduction} itself against a case
with an independently known answer: applied to $q^2>4pr$ on $[-1,1]^3$ (also a
cone, the quadratic discriminant being homogeneous), the three faces give
$S_p=S_r=13/6$, $S_q=5/2+\ln2$, and $\tfrac13\cdot2(S_p+S_q+S_r)$ works out to
$\tfrac{41}{9}+\tfrac23\ln2$ -- exactly eight times the known full-quadratic
volume $\tfrac{41}{72}+\tfrac{\ln2}{12}$. The reduction is an identity, not an
approximation.
\end{remark}

\subsection{Computing $S_b$}

Substitute $s=\sqrt{1-3ac}$ (so $c=(1-s^2)/(3a)$). The admissible-$d$ band on
this face is $[-K_-/u,\,K_+/u]$ with
\[
  K_+=(s-1)^2(2s+1),\qquad K_-=(s+1)^2(2s-1),\qquad u=27a^2,\qquad K_++K_-=4s^3,
\]
on the domain $\{s\in(0,2),\ a\in[|s^2-1|/3,1]\}$. Integrating over $a$
\emph{first}, at fixed $s$, is elementary, since $a$ enters only through
$u=27a^2$; two lemmas make it precise.

\begin{lemma}[Never clips above] \label{lem:L1}
$d_{\mathrm{hi}}\le1$ throughout: $\alpha_+:=\sqrt{K_+/27}<a_0:=|s^2-1|/3$
reduces to $2s+1<3(s+1)^2$, true for all $s$. This is the exact analogue of
Lemma~\ref{lem:never-clipped}.
\end{lemma}

\begin{lemma}[Clips below exactly for $s\in(2/3,2)$] \label{lem:L2}
$\alpha_-:=\sqrt{K_-/27}>a_0$ reduces to $(3s-2)(s-2)<0$; and $\alpha_-<1$
reduces to $(s+1)^2(2s-1)<27$, i.e.\ $s<2$, with equality precisely at
$s=2$.
\end{lemma}

With Lemmas~\ref{lem:L1}--\ref{lem:L2}, $F(s):=\int_{a_0}^1 L(a,s)/a\,da$ (the
$a$-integral of the band length) is elementary -- the term
$-K_+/(2K_-)+2s^3/K_-$ collapses to the constant $1/2$ via $K_++K_-=4s^3$ --
giving the piecewise closed form
\[
  F(s)=
  \begin{cases}
    \dfrac{2s^3}{27}\left(\dfrac{1}{a_0^2}-1\right), & s\le2/3, \\[2ex]
    \dfrac{K_+}{54a_0^2}+\dfrac12+\ln\dfrac{\alpha_-}{a_0}-\dfrac{2s^3}{27}, & 2/3<s<2,
  \end{cases}
\]
continuous at $s=2/3$ (both branches equal $11264/18225$ there). Then
\[
  S_b = \frac43\int_0^2 s\,F(s)\,ds = \frac{1454}{405}-\frac56\ln3
      = 2.674613216233365\ldots,
\]
and, substituting into \eqref{eq:cone-reduction} with
$V(1)=\tfrac{766}{1215}+\tfrac{\ln3}{6}$,
\[
  \Prb(\text{non-monic}) = \frac{1}{16}\left(\frac{766}{1215}+\frac{\ln3}{6}+\frac{1454}{405}-\frac56\ln3\right)
  = \frac{641}{2430}-\frac{\ln3}{24}. \qedhere
\]
(Confirmed to hold exactly by direct symbolic simplification in \texttt{sympy}.)

\subsection{Independent verification}

Two of the checks below are fully independent of each other and of the
derivation above; the rest are internal consistency checks of the
computation.

\begin{enumerate}
\item \textbf{Exact symbolic evaluation} (as above) -- \texttt{sympy} confirms
  $\Prb-(\tfrac{641}{2430}-\tfrac{\ln3}{24})=0$ identically.
\item \textbf{An independent numerical route, sharing no step with the
  derivation above}: fixing the leading coefficient $a$ and integrating $a$
  out \emph{last} (rather than first, as in \eqref{eq:cone-reduction}) --
  effectively a direct high-precision evaluation of
  \eqref{eq:naive-integral}, with the $a\to0$ endpoint singularity and the
  migrating clipping breakpoints of the inner integral carefully handled --
  reproduces $0.21801049620261477102\ldots$, agreeing with the exact value to
  \emph{nineteen} significant digits (Gauss--Legendre and tanh--sinh
  quadrature, several degrees each, spread $1.6\times10^{-18}$).
\item \textbf{Blind PSLQ}: given only the $22$-digit numeric value of $S_b$,
  a PSLQ integer-relation search (no closed form supplied) recovers
  $S_b=\tfrac{1454}{405}-\tfrac56\ln3$ unprompted.
\item \textbf{Monte Carlo}, $4\times10^{10}$ samples on the raw sign of
  \eqref{eq:disc4} (using no theory from this paper): $0.218008985\pm2.1\times10^{-6}$,
  $-0.73\sigma$ from the exact value.
\item Double-precision brute-force quadrature (no special handling of any
  kind): $0.2180105004$, $4.2\times10^{-9}$ from the exact value -- at that
  method's own precision limit.
\end{enumerate}

\begin{remark}
This integrand carries three numerical hazards worth recording, all
concentrated at small leading coefficient $a$ and all invisible to a
quadrature routine's own convergence estimate: (i) the conditional
probability given the leading coefficient has a \emph{square-root} endpoint
singularity as $a\to0$, $p(a)=p(0)-\tfrac23\sqrt a+O(a)$ with
$p(0)=\tfrac12+\tfrac{5}{72}+\tfrac{\ln2}{12}$ (the random-quadratic limit,
since as $a\to0$ the third root escapes to $-b/a$), on which a fixed-degree
Gauss rule loses roughly seven digits while reporting no warning; (ii) the
clipping/kink structure in the inner variable migrates to scale
$\sim\sqrt{3a}$ as $a\to0$, so any fixed panel layout eventually walks past
it; and (iii) floating-point arithmetic loses the clipping breakpoints
altogether for small $a$, since the relevant breakpoint cubic has a
near-double root whose branches separate by only $O(\sqrt a)$. All three had
to be identified and handled to reach the nineteen-digit confirmation above.
\end{remark}

\subsection{Relation to the \texttt{dxdy.ru} thread}
The forum thread's stated intermediate volume,
$V_0=\tfrac{766}{1215}+\tfrac{\ln3}{6}$, agrees with $V(1)=8\cdot$
Theorem~\ref{thm:sym-cube}, and its final claimed value,
$\tfrac{641}{2430}-\tfrac{\ln3}{24}$, is confirmed exactly by
Theorem~\ref{thm:nonmonic}. Its remaining steps therefore do reach the right
answer, though nothing in the present paper validates the specific route the
thread took to get there; the derivation above is independent.

\section{Formal verification in Lean} \label{sec:lean}

All three theorems have been formalized and \emph{completely, machine-checked
proved} in Lean~4 against Mathlib (toolchain \texttt{v4.33.0}), with
\textbf{zero} uses of \texttt{sorry} anywhere in the development and no
unexpected axioms: \texttt{\#print axioms} on every headline theorem returns
exactly \texttt{[propext, Classical.choice, Quot.sound]}, the three standard
axioms used throughout Mathlib itself. The development is $3780$ lines across
seven files (\texttt{Basic.lean}, \texttt{Integrals.lean},
\texttt{Theorem1.lean}, \texttt{Theorem2.lean},
\texttt{Theorem3Statement.lean}, \texttt{Theorem3Proof.lean},
\texttt{DiscriminantRootCount.lean}), all in the accompanying repository, and
rebuilds from a warm Mathlib cache in under a minute.

\subsection*{The discriminant--root-count bridge}
Everywhere above, ``has three real roots'' is proved as a statement about the
\emph{sign of $\disc$}. Mathlib has \texttt{Cubic.discr} and a criterion for
$\texttt{discr}\ne0$ implying distinct roots over a splitting field, but
nothing connecting the \emph{sign} of a real cubic's discriminant to how many
of its roots are real -- the fact stated without proof in
Section~\ref{sec:prelim} above, and standard since at least
Cardano. \texttt{DiscriminantRootCount.lean} proves it from scratch and with
no appeal to calculus: for $p=b^2-3ac>0$, evaluating $27a^2f$ at the two
points where $3at+b=\mp\sqrt p$ gives values whose \emph{product} is
$-27a^2\disc$ and whose \emph{difference} is $4p^{3/2}$, and three
applications of the intermediate value theorem to the resulting sign pattern
produce three real roots directly, with no derivative or critical point
appearing anywhere. Theorems~\ref{thm:sym-cube} and \ref{thm:unit-cube} are
additionally restated and proved as \texttt{theorem1\_root\_count} and
\texttt{theorem2\_root\_count}, in terms of \texttt{HasThreeDistinctRealRoots}
directly, with no discriminant in the statement at all.

\subsection*{Two presentational simplifications found while formalizing}
The informal proof of Theorems~\ref{thm:sym-cube}--\ref{thm:unit-cube} above
derives the band via the critical points $x_\pm$ of $g$. In Lean, it proved
cheaper to skip $x_\pm$ entirely and use the single \texttt{ring}-provable
identity
\[
  -27\,\disc_3(a,b,c) = (27c-9ab+2a^3)^2 - 4(a^2-3b)^3,
\]
which yields both branches of \eqref{eq:disc-factor} at once, and to prove
Lemma~\ref{lem:never-clipped} from the closed forms
$27c_{\mathrm{hi}}=(a-s)^2(a+2s)$, $27c_{\mathrm{lo}}=(a+s)^2(a-2s)$ directly,
reducing the whole lemma to a single scalar inequality on $a\in[-1,1]$,
$s\in[0,2]$ that \texttt{nlinarith} discharges automatically with product
hints.

More strikingly, the cone step of \S\ref{sec:nonmonic}
(\eqref{eq:cone-reduction}) needs \emph{no divergence theorem at all}. The
Lean proof of \texttt{volume\_conePiece1} establishes, for any measurable
cone $R\subseteq\RR^4$, that the piece of $[-1,1]^4$ where coordinate $1$
attains $\max_j|x_j|$ has volume $\tfrac14\vol_3(\text{face})$, by slicing at
$x_1=t$, using homogeneity of $R$ to identify the slice with $t\cdot(\text{face})$,
and integrating $\int_0^1t^3\,dt=\tfrac14$. This is Fubini plus the elementary
scaling law $\vol_3(t\cdot E)=t^3\vol_3(E)$: no boundary integral appears
anywhere, so the cusps of $\{\disc=0\}$ never have to be dealt with, and no
smoothness of $\partial R$ is used or needed. The seven remaining pieces
follow from the first by measure-preserving coordinate maps built from the
two \texttt{ring}-provable symmetries of $\disc$. \emph{Both} simplifications
are the same mathematics as the informal derivation, just in a form that
turned out to need less machinery than either the critical-point picture or
the classical divergence theorem statement suggested -- and neither changes
any of the numbers above.

\begin{remark}
Formalizing surfaced one genuine error and confirms none of the paper's
numerical claims were wrong. An early Lean statement of the inner integral
$F(s)$ (\eqref{eq:naive-integral}'s companion, defined for $s\in(0,2)$) was
false exactly at $s=1$, where $a_0=0$ and the integral genuinely diverges;
both sides collapsed to unrelated junk values under Lean's total-function
conventions ($0/0=0$, $\log0=0$) rather than raising an error, and only
attempting the proof -- not merely writing the statement -- caught it. The
fix is an $s\ne1$ hypothesis, harmless since $\{1\}$ is null. We record this
because it is a direct illustration of why every numerical claim in this
paper rests on multiple independent computations rather than on any one of
them: an unchecked claim, even one that type-checks, is not yet a fact.
\end{remark}

\section{Novelty and prior art} \label{sec:novelty}

An eight-agent parallel literature sweep (arXiv, journal databases, Math
StackExchange and MathOverflow's own search indices, OEIS, and general web
search), followed by a targeted search on the exact constants derived here,
found:

\begin{itemize}
\item \textbf{Theorem~\ref{thm:sym-cube}}: no prior statement of this
  probability anywhere searchable (OEIS direct digit-sequence query: no
  match; Math StackExchange/MathOverflow site search for the monic cubic:
  zero threads; web search for the exact fraction and decimal: empty). Two
  caveats. First, the \emph{unnormalized} volume $V_0=\tfrac{766}{1215}+\tfrac{\ln3}{6}$
  (i.e.\ eight times our value) appears, uninterpreted as a probability, as an
  intermediate step of the \texttt{dxdy.ru} thread discussed above -- an
  unrefereed, partly AI-assisted forum derivation of the \emph{non-monic}
  cubic, dated June 2026, not indexed by any search engine at the time of our
  search. Second, Li~\cite{li1988} claims an exact expression for the cubic
  family $x^3+3ax^2+3bx+2c$ with $(a,b,c)$ uniform on symmetric intervals
  $[-h,h]\times[-k,k]\times[-l,l]$; at $h=k=1/3$, $l=1/2$ this specializes to
  exactly our model. The paper itself is paywalled and, per its recovered
  \emph{Mathematical Reviews} notice \cite{li1988-mr}, no numerical or
  symbolic value from it has ever been reproduced by either of its two
  citations; whether it correctly reduces to Theorem~\ref{thm:sym-cube}
  cannot be checked without the original text.
\item \textbf{Theorem~\ref{thm:unit-cube}}: likely new without caveat -- Li's
  furnished special cases are all symmetric intervals, and the value does not
  appear in the \texttt{dxdy.ru} thread or anywhere else searched.
\item \textbf{Theorem~\ref{thm:nonmonic}}: the \emph{value} is not new (see
  above); the \emph{proof} and its independent verification are.
\end{itemize}

We regard the Li~\cite{li1988} question as the one genuinely open item in this
account.

\section{Open problems}

The reduction techniques used here appear to generalize.

\begin{enumerate}
\item \textbf{Monic quartic.} The quartic discriminant is homogeneous of
  degree~$6$ in the \emph{five} coefficients $a,b,c,d,e$ of
  $ax^4+bx^3+cx^2+dx+e$, so the divergence-theorem argument of
  \S\ref{sec:nonmonic} does apply -- but to the \emph{non-monic} quartic on
  $[-1,1]^5$, reducing that all-real-roots volume to a sum over $2\times5=10$
  three-dimensional face volumes. One of those faces (the one with the
  leading coefficient pinned to $1$) is exactly the \emph{monic} quartic
  probability we would actually like; the reduction therefore needs it as an
  \emph{input}, not a consequence, and does not apply directly to the monic
  problem stated here. Fixing the leading coefficient at $1$ breaks the
  homogeneity the cone argument relies on, so the monic quartic likely needs
  a direct generalization of the band argument of \S3 instead -- complicated
  by a quartic's derivative being itself a cubic, with up to three real
  critical points rather than one band. Current Monte Carlo estimate:
  $0.0054749\pm9.0\times10^{-6}$; under active investigation at the time of
  writing.
\item \textbf{Gaussian coefficients.} For the monic cubic with $(a,b,c)$
  i.i.d.\ standard normal, there is no cube to bound the band, so
  Lemma~\ref{lem:never-clipped} does not apply in any form; the natural
  object is $\mathbb{E}[\Phi(c_{\mathrm{hi}})-\Phi(c_{\mathrm{lo}})]$, a smooth
  two-dimensional integral with no case splits, but a priori no reason to
  expect an elementary closed form. Current Monte Carlo estimate:
  $0.169962\pm4.2\times10^{-5}$; under active investigation at the time of
  writing.
\item \textbf{General degree $n$.} The probability of no real roots for a
  degree-$n$ Kac polynomial decays as $n^{-3/4}$ exactly
  (Poplavskyi--Schehr~\cite{poplavskyi-schehr}); the corresponding rate
  constant for \emph{all} $n$ roots real, under speed-$n^2$ large-deviation
  scaling, is not known in closed form for uniform (or Kac-Gaussian)
  coefficients.
\item \textbf{The parent matrix problem.} For an $n\times n$ matrix with
  i.i.d.\ uniform entries, the all-real-eigenvalue probability remains open
  for $n\ge3$; the striking empirical near-coincidence
  $\Prb_n(\text{uniform})\approx2^{-(n-1)(n-2)/4}$ (Edelman's Gaussian
  probability at rank $n-1$), first observed numerically on Math
  StackExchange, has no known explanation.
\end{enumerate}

\appendix
\section{Numerical verification summary}

\begin{center}
\renewcommand{\arraystretch}{1.3}
\begin{tabular}{lccc}
\hline
 & Theorem~\ref{thm:sym-cube} & Theorem~\ref{thm:unit-cube} & Theorem~\ref{thm:nonmonic} \\
\hline
Symbolic (\texttt{sympy}, exact) & PASS & PASS & PASS ($\Prb-\text{cand.}=0$) \\
Independent numerical route & $\sim10^{-16}$ & $\sim10^{-13}$ & 19 sig.\ digits ($7.0\times10^{-20}$) \\
Blind PSLQ / \texttt{identify} & --- & --- & recovers exact form \\
Monte Carlo & $2\times10^8$: $+1.3\sigma$ & $2\times10^8$: $-0.5\sigma$ & $4\times10^{10}$: $-0.73\sigma$ \\
Lean~4 / Mathlib & \textbf{0 \texttt{sorry}} & \textbf{0 \texttt{sorry}} & \textbf{0 \texttt{sorry}} \\
\hline
\end{tabular}
\end{center}

\section*{Data and code availability}
All source code (Python: \texttt{sympy}/\texttt{mpmath}/\texttt{numpy}/\texttt{scipy};
Lean~4: Mathlib \texttt{v4.33.0}), raw numerical results, the literature-sweep
records underlying \S\ref{sec:novelty}, and an interactive visual explainer of
the monic-cubic case are available in the accompanying repository.

\begin{thebibliography}{9}

\bibitem{li1988}
Hung C.~Li, \emph{The exact probability that the roots of quadratic, cubic,
and quartic equations are all real if the equation coefficients are random},
Comm.\ Statist.\ Theory Methods \textbf{17} (1988), no.~2, 395--409.

\bibitem{li1988-mr}
G.~Samal, review of \cite{li1988}, Mathematical Reviews MR89j:60069 (1989).

\bibitem{haldar-chakraborty}
S.~Haldar and S.~Chakraborty, \emph{On the Probability of Real Roots in a
Quadratic Equation with Coefficients as i.i.d.\ $U(-\theta,\theta)$
Variates}, J.\ Appl.\ Math.\ Comput.\ \textbf{5} (2021), no.~1, 48--55.

\bibitem{hetzel-liew-morrison}
A.~J.~Hetzel, J.~S.~Liew, and K.~E.~Morrison, \emph{The Probability that a
Matrix of Integers Is Diagonalizable}, Amer.\ Math.\ Monthly \textbf{114}
(2007), 491--499.

\bibitem{edelman1997}
A.~Edelman, \emph{The Probability that a Random Real Gaussian Matrix has $k$
Real Eigenvalues, Related Distributions, and the Circular Law}, J.\
Multivariate Anal.\ \textbf{60} (1997), no.~2, 203--232.

\bibitem{mse-eigenvalue-question}
\emph{Probability for an $n\times n$ matrix to have only real eigenvalues},
Mathematics Stack Exchange, question 3770846, asked 2020-07-27.

\bibitem{poplavskyi-schehr}
M.~Poplavskyi and G.~Schehr, \emph{Exact Persistence Exponent for the
2d-Diffusion Equation and Related Kac Polynomials}, Phys.\ Rev.\ Lett.\
\textbf{121} (2018), 150601.

\end{thebibliography}

\end{document}
